% Unidad IV — 4.2 Foro Dinámica de Liderazgo (5 %); numeración local: 9.
% Técnica: foro. Nivel: analizar, fundamentar y dialogar.
% Organización: problema -> postura y tres respuestas -> interacción -> conclusión.
% Estado: preparación local; no acredita publicación ni interacción en Moodle.
\documentclass[11pt,letterpaper,oneside]{article}
\usepackage[margin=2.35cm]{geometry}
\usepackage[utf8]{inputenc}
\usepackage[T1]{fontenc}
\usepackage[spanish,es-tabla]{babel}
\usepackage[scaled=.96]{helvet}
\renewcommand{\familydefault}{\sfdefault}
\usepackage{graphicx}
\usepackage{xcolor}
\usepackage{microtype}
\usepackage{enumitem}
\usepackage[most]{tcolorbox}
\usepackage[authoryear,round]{natbib}
\usepackage{xurl}
\usepackage{hyperref}
\usepackage{attachfile}
\definecolor{itescaRedDark}{HTML}{8F1117}
\definecolor{itescaRed}{HTML}{D70F18}
\definecolor{itescaPaper}{HTML}{F7F7F7}
\setlength{\parindent}{0pt}
\setlength{\parskip}{.6em}
\setlist[itemize]{leftmargin=*,itemsep=.25em,topsep=.3em}
\newcommand{\itescadocumenttitle}{Dinámica de liderazgo}
\newcommand{\itescadocumentsubtitle}{4.2 Foro · Unidad IV · 5\%}
\newcommand{\itescadocumentsubject}{Actividad 9 - Fundamentos de Gestión Administrativa}
\hypersetup{pdftitle={\itescadocumenttitle},pdfsubject={\itescadocumentsubject},
  pdfauthor={Martin Jonathan de la Cruz},colorlinks=true,
  linkcolor=itescaRedDark,citecolor=itescaRedDark,urlcolor=itescaRedDark}
\newtcolorbox{forobox}[1]{enhanced,breakable,colback=itescaPaper,
  colframe=itescaRedDark,boxrule=.5pt,arc=1mm,
  title={#1},fonttitle=\bfseries,left=3mm,right=3mm,top=2mm,bottom=2mm}
\bibliographystyle{plainnat}

\begin{document}
\begin{center}
  \includegraphics[width=4cm]{ITESCA/assets-itesca/web/logo-itesca-contorno.png}\\[.4em]
  {\small Instituto Tecnológico Superior de Cajeme · Maestría en Gestión Administrativa}\\[.8em]
  {\LARGE\bfseries\color{itescaRedDark}\itescadocumenttitle}\\[.3em]
  {\large\itescadocumentsubtitle}\\[.5em]
  Fundamentos de Gestión Administrativa · Curso GGFG02\\
  Martin Jonathan de la Cruz\\[.4em]
  {\small Preparación: 12 de septiembre de 2026 · Vencimiento indicado: 7 de septiembre de 2026, 23:59}
\end{center}

\section{Introducción}
Dirigir un equipo diverso exige combinar resultados, aprendizaje y trato justo. La autoridad formal permite asignar tareas, pero no garantiza compromiso ni buenas decisiones. El problema consiste en elegir conductas de liderazgo acordes con la tarea, el riesgo y las capacidades de las personas, sin convertir la adaptación en arbitrariedad. La clasificación divulgativa de Conduce Tu Empresa distingue estilos como el democrático, transformacional y transaccional; conviene contrastarla con la literatura y no tratar sus etiquetas como categorías excluyentes \citep{conduceEstilosLiderazgo,openstaxContingencia2019}.

\section{Dirigir con propósito, participación y adaptación}
\subsection{Aportación inicial}
\textattachfile[description={Aportación inicial en texto copiable},mimetype=text/plain]{ITESCA/maestria-en-gestion-administrativa/fundamentos-de-gestion-administrativa-mga/foro-participacion-Actividad-9.txt}{\textcolor{itescaRedDark}{\textbf{Abrir aportación en texto copiable}}}

\begin{forobox}{Liderar con flexibilidad, sin renunciar a la responsabilidad}
Hola, compañeras y compañeros. Considero que un liderazgo eficaz no consiste en imponer siempre el mismo estilo, sino en combinar propósito, participación y acompañamiento según la situación, manteniendo criterios éticos constantes.

\textbf{1. En un puesto de nivel medio o alto, ¿qué estilo o estilos de liderazgo utilizaría y por qué?}

Utilizaría principalmente un liderazgo \textbf{transformacional y democrático o participativo}, con adaptación a la tarea y a las capacidades del equipo. El primero permite comunicar un propósito compartido, estimular mejoras y desarrollar a las personas; el segundo incorpora el conocimiento de quienes realizan el trabajo antes de decidir. Participar no significa someter toda decisión a votación ni diluir la responsabilidad del directivo \citep{openstaxTransformacional2019,openstaxContingencia2019}.

En un supuesto de digitalización del área administrativa, explicaría qué problema se busca resolver, consultaría a quienes conocen el proceso y delegaría decisiones operativas a quienes demuestren competencia. A una persona que aprende una tarea nueva le ofrecería instrucciones y acompañamiento, sin asumir su capacidad por la edad. Esto coincide con la teoría trayectoria--meta: aclarar objetivos, retirar obstáculos y ajustar la conducta directiva a las necesidades de la situación \citep{openstaxContingencia2019}.

Complementaría esos estilos con acuerdos transaccionales claros sobre calidad, plazos y reconocimiento. Ante una contingencia que requiera respuesta inmediata, asumiría una conducción más directiva, limitada al riesgo y sujeta a revisión posterior. No adoptaría el autoritarismo permanente ni confundiría delegación con ausencia de seguimiento: la flexibilidad debe cambiar el apoyo y la participación, no la dignidad ni las reglas de justicia.

\textbf{2. ¿Qué cualidades se requieren para dirigir exitosamente una organización o equipo?}

Priorizaría cinco cualidades conectadas con conductas observables:
\begin{itemize}
  \item \textbf{Integridad y responsabilidad:} cumplir acuerdos, explicar criterios y reconocer errores para sostener la confianza.
  \item \textbf{Empatía y escucha:} identificar dificultades sin descalificar, resolver conflictos y comprender necesidades antes de exigir resultados.
  \item \textbf{Visión y comunicación clara:} convertir el propósito en prioridades, responsabilidades y expectativas comprensibles.
  \item \textbf{Juicio analítico:} contrastar datos, experiencia y riesgos, evitando decidir solo por intuición o popularidad.
  \item \textbf{Adaptabilidad y capacidad de desarrollar a otros:} ofrecer retroalimentación, formación y autonomía gradual, sin hacer al equipo dependiente del jefe.
\end{itemize}
Estas cualidades no garantizan por sí solas el éxito: deben traducirse en prácticas y condiciones de trabajo. La consideración individual y la estimulación intelectual sustentan el liderazgo transformacional, mientras que el apoyo y la claridad de funciones son centrales en la trayectoria--meta \citep{openstaxTransformacional2019,openstaxContingencia2019}. Evaluaría su efecto mediante calidad del servicio, cumplimiento de acuerdos y percepción de trato justo, no solamente con resultados financieros.

\textbf{3. Ante las brechas generacionales y las megatendencias, ¿qué cambios haría en los procesos de dirección?}

Sánchez identifica diferencias de comunicación, motivación y formación, y recomienda colaboración intergeneracional y desarrollo personalizado \citep{sanchezRaona2023}. Considero que esto exige diagnosticar necesidades individuales: pertenecer a una generación no determina la habilidad digital, el compromiso ni las preferencias de una persona.

Además, el Foro Económico Mundial identifica cambio tecnológico, transición verde, cambios demográficos e incertidumbre económica entre las tendencias que transformarán el trabajo. Son expectativas de empleadores para 2025--2030, no resultados garantizados para cada organización \citep{wefFutureJobs2025}. A partir de ellas propondría:
\begin{itemize}
  \item \textbf{Comunicación accesible y participación:} acordar canales, documentar decisiones y combinar conversaciones con mensajes asincrónicos para evitar exclusiones.
  \item \textbf{Aprendizaje continuo y mentoría recíproca:} intercambiar experiencia de procesos y competencias digitales según habilidades comprobadas, no estereotipos de edad.
  \item \textbf{Delegación con resultados verificables:} acordar entregables, recursos y revisiones; permitir flexibilidad cuando la tarea lo admita y evitar medir compromiso por presencia permanente.
  \item \textbf{Dirección responsable del cambio:} probar la automatización en pequeña escala, proteger datos y mantener revisión humana; incorporar ahorro de recursos, bienestar y planes de contingencia a los objetivos del equipo.
\end{itemize}
En el supuesto de digitalización, la gerencia asignaría recursos y límites; cada responsable de área acordaría metas y formación con su equipo. Compararía mensualmente errores, tiempos de atención y aprendizaje aplicado con una línea base. Si mejora la velocidad, pero aumentan errores o sobrecarga, corregiría el proceso antes de ampliarlo.

En síntesis, dirigir exitosamente es desarrollar autonomía con responsabilidad. ¿Cómo equilibrarían ustedes la participación del equipo con una decisión urgente que no puede esperar al consenso?
\end{forobox}

\subsection{Dos réplicas modelo para adaptar}
\textbf{No son respuestas publicadas ni acreditan interacción.} Antes de utilizarlas, seleccionar dos aportaciones reales, mencionar con fidelidad una idea de cada una y ajustar el argumento. No se atribuyen opiniones a compañeros inexistentes.

\begin{forobox}{Modelo 1 · Si una aportación defiende la participación}
El liderazgo participativo puede mejorar la decisión cuando aprovecha la experiencia del equipo, aunque consultar no elimina la responsabilidad de quien dirige. Añadiría un criterio: ajustar el alcance de la consulta al riesgo, al tiempo disponible y al conocimiento de los participantes, como permite razonar la teoría trayectoria--meta \citep{openstaxContingencia2019}. En una urgencia propondría escuchar a las personas clave, decidir y explicar después los criterios. ¿Qué decisiones reservarías al responsable y cuáles delegarías para evitar que la participación retrase innecesariamente el trabajo?
\end{forobox}

\begin{forobox}{Modelo 2 · Si una aportación aborda generaciones o tecnología}
La colaboración entre edades puede enriquecer el equipo, pero conviene evitar que las etiquetas generacionales sustituyan la evaluación individual. La formación personalizada y la colaboración intergeneracional recomendadas por Sánchez permiten plantear mentoría recíproca basada en conocimientos reales \citep{sanchezRaona2023}. Propondría parejas que intercambien experiencia de procesos y habilidades digitales, con una mejora concreta como resultado. ¿Cómo comprobarías que ambas personas aprenden y que el programa no refuerza estereotipos ni asigna siempre los mismos papeles?
\end{forobox}

\clearpage
\section{Conclusiones}
El criterio de dirección más sólido combina propósito transformacional, participación y adaptación de las conductas de apoyo y supervisión. Su límite es ético: ni la urgencia justifica el trato arbitrario ni la flexibilidad permite abandonar la responsabilidad. Las diferencias entre personas y los cambios tecnológicos demandan aprendizaje, comunicación accesible y evaluación conjunta de resultados y bienestar. Mi postura es delegar con recursos, criterios y seguimiento, y ajustar la intervención según evidencia, no según estereotipos. Esta conclusión expresa una posición previa al intercambio; no una síntesis de conversaciones que aún no se han verificado.\footnote{Se utilizó inteligencia artificial como apoyo para organizar y redactar el borrador. Corresponde al estudiante revisar las fuentes, contrastar el video completo, asumir o modificar la postura y adaptar las réplicas antes de publicar.}

\bibliography{fundamentos-de-gestion-administrativa}
\end{document}
